% Môn: Vật lý
% Lớp: 10
% Chương: Chương I. Mở đầu
% Bài: L10C1 Bài 2. Các quy tắc an toàn trong phòng thực hành Vật lí · Nhận biết và giải thích biển báo, kí hiệu cảnh báo
% Số câu: 185
% App đọc file này trực tiếp — sửa rồi Commit trên GitHub, bấm Đồng bộ GitHub .tex

% L10C1 Bài 2. Các quy tắc an toàn trong phòng thực hành Vật lí



% ===== Câu 97 =====
% ID: L10C1B2-25-DS
% Mức: NB
\dangbt{Nhận biết và giải thích biển báo, kí hiệu cảnh báo}
\begin{ex}
% ID: L10C1B2-25-DS
% Mức: NB
Trong tình huống khi khu vực làm việc bị ướt đối với biển báo và kí hiệu cảnh báo, hãy đánh giá các phát biểu sau.
\choiceTF
{\True Cần quan sát hình dạng, màu sắc, biểu tượng rồi thực hiện đúng chỉ dẫn.}
{Có thể bỏ qua biển cảnh báo vì cho rằng chỉ mang tính tham khảo nếu thao tác diễn ra nhanh.}
{\True Phải báo cho giáo viên hoặc người phụ trách khi phát hiện nguy cơ vượt khả năng kiểm soát.}
{\True Chỉ tiếp tục thí nghiệm sau khi nguyên nhân nguy hiểm đã được loại bỏ hoặc kiểm soát.}
\loigiai{
\begin{itemchoice}
\item (Đúng) Cần quan sát hình dạng, màu sắc, biểu tượng rồi thực hiện đúng chỉ dẫn. (Vì): vì đây là biện pháp an toàn của dạng biển báo và kí hiệu cảnh báo; b sai vì hành vi “bỏ qua biển cảnh báo vì cho rằng chỉ mang tính tham khảo” vẫn nguy hiểm; c và d đúng do sự cố phải được báo cáo, cô lập và kiểm soát trước khi tiếp tục
\item (Sai) Có thể bỏ qua biển cảnh báo vì cho rằng chỉ mang tính tham khảo nếu thao tác diễn ra nhanh.
\item (Đúng) Phải báo cho giáo viên hoặc người phụ trách khi phát hiện nguy cơ vượt khả năng kiểm soát.
\item (Đúng) Chỉ tiếp tục thí nghiệm sau khi nguyên nhân nguy hiểm đã được loại bỏ hoặc kiểm soát.
\end{itemchoice}
}
\end{ex}

% ===== Câu 99 =====
% ID: L10C1B2-26-DS
% Mức: TH
\begin{ex}
% ID: L10C1B2-26-DS
% Mức: TH
Trong tình huống khi cần rời bàn thí nghiệm đối với biển báo và kí hiệu cảnh báo, hãy đánh giá các phát biểu sau.
\choiceTF
{\True Cần quan sát hình dạng, màu sắc, biểu tượng rồi thực hiện đúng chỉ dẫn.}
{Có thể bỏ qua biển cảnh báo vì cho rằng chỉ mang tính tham khảo nếu thao tác diễn ra nhanh.}
{\True Phải báo cho giáo viên hoặc người phụ trách khi phát hiện nguy cơ vượt khả năng kiểm soát.}
{\True Chỉ tiếp tục thí nghiệm sau khi nguyên nhân nguy hiểm đã được loại bỏ hoặc kiểm soát.}
\loigiai{
\begin{itemchoice}
\item (Đúng) Cần quan sát hình dạng, màu sắc, biểu tượng rồi thực hiện đúng chỉ dẫn. (Vì): vì đây là biện pháp an toàn của dạng biển báo và kí hiệu cảnh báo; b sai vì hành vi “bỏ qua biển cảnh báo vì cho rằng chỉ mang tính tham khảo” vẫn nguy hiểm; c và d đúng do sự cố phải được báo cáo, cô lập và kiểm soát trước khi tiếp tục
\item (Sai) Có thể bỏ qua biển cảnh báo vì cho rằng chỉ mang tính tham khảo nếu thao tác diễn ra nhanh.
\item (Đúng) Phải báo cho giáo viên hoặc người phụ trách khi phát hiện nguy cơ vượt khả năng kiểm soát.
\item (Đúng) Chỉ tiếp tục thí nghiệm sau khi nguyên nhân nguy hiểm đã được loại bỏ hoặc kiểm soát.
\end{itemchoice}
}
\end{ex}

% ===== Câu 100 =====
% ID: L10C1B2-26-TL
% Mức: TH
\begin{bt}
% ID: L10C1B2-26-TL
% Mức: TH
Phân tích hành vi “bỏ qua biển cảnh báo vì cho rằng chỉ mang tính tham khảo” và đề xuất cách thay thế an toàn.
\loigiai{
Hành vi “bỏ qua biển cảnh báo vì cho rằng chỉ mang tính tham khảo” có thể gây tai nạn vì làm mất kiểm soát nguồn nguy hiểm. Cần dừng thao tác, cô lập nguy cơ và thay bằng biện pháp: quan sát hình dạng, màu sắc, biểu tượng rồi thực hiện đúng chỉ dẫn.
}
\end{bt}

% ===== Câu 101 =====
% ID: L10C1B2-27-TL
% Mức: VD
\begin{bt}
% ID: L10C1B2-27-TL
% Mức: VD
Xây dựng một bảng kiểm ngắn gồm ít nhất bốn mục cho tình huống khi cần rời bàn thí nghiệm liên quan đến biển báo và kí hiệu cảnh báo.
\loigiai{
Bảng kiểm gồm: đọc hướng dẫn; nhận diện mối nguy; kiểm tra dụng cụ và bảo hộ; thực hiện quan sát hình dạng, màu sắc, biểu tượng rồi thực hiện đúng chỉ dẫn; theo dõi bất thường; thu dọn và báo cáo sau thí nghiệm.
}
\end{bt}

% ===== Câu 103 =====
% ID: L10C1B2-28-TL
% Mức: VD
\begin{bt}
% ID: L10C1B2-28-TL
% Mức: VD
Một học sinh đã bỏ qua biển cảnh báo vì cho rằng chỉ mang tính tham khảo. Hãy nêu trình tự xử lí ngay sau khi phát hiện hành vi này.
\loigiai{
Trình tự: yêu cầu dừng ngay; không tiếp cận nếu chưa an toàn; cô lập nguồn nguy hiểm; báo giáo viên; sơ cứu hoặc xử lí theo hướng dẫn; chỉ hoạt động lại sau khi được xác nhận an toàn.
}
\end{bt}

% ===== Câu 107 =====
% ID: L10C1B2-30-TL
% Mức: VDC
\begin{bt}
% ID: L10C1B2-30-TL
% Mức: VDC
Đề xuất cách tổ chức nhóm học sinh để thực hiện nhiệm vụ về biển báo và kí hiệu cảnh báo an toàn và có thể kiểm tra được.
\loigiai{
Phân công trưởng nhóm kiểm tra điều kiện, người thao tác, người quan sát an toàn và người ghi chép. Dùng bảng kiểm, xác nhận trước mỗi bước, quy định tín hiệu dừng và báo ngay bất thường.
}
\end{bt}

% ===== Câu 152 =====
% ID: L10C1B2-55-TN
% Mức: NB
\begin{ex}
% ID: L10C1B2-55-TN
% Mức: NB
Trước khi bắt đầu thí nghiệm liên quan đến biển báo và kí hiệu cảnh báo, hành động nào an toàn nhất?
\choice
{\True quan sát hình dạng, màu sắc, biểu tượng rồi thực hiện đúng chỉ dẫn}
{bỏ qua biển cảnh báo vì cho rằng chỉ mang tính tham khảo}
{tiếp tục thao tác rồi báo cáo sau}
{nhờ bạn khác làm thay mà không kiểm tra điều kiện an toàn}
\loigiai{
Hành động đúng là “quan sát hình dạng, màu sắc, biểu tượng rồi thực hiện đúng chỉ dẫn” vì giúp nhận biết nhanh loại nguy hiểm, điều cấm và hành động bắt buộc. Chọn A.
}
\end{ex}

% ===== Câu 153 =====
% ID: L10C1B2-56-TN
% Mức: NB
\begin{ex}
% ID: L10C1B2-56-TN
% Mức: NB
Khi phát hiện thiết bị có dấu hiệu bất thường liên quan đến biển báo và kí hiệu cảnh báo, hành động nào an toàn nhất?
\choice
{bỏ qua biển cảnh báo vì cho rằng chỉ mang tính tham khảo}
{tiếp tục thao tác rồi báo cáo sau}
{nhờ bạn khác làm thay mà không kiểm tra điều kiện an toàn}
{\True quan sát hình dạng, màu sắc, biểu tượng rồi thực hiện đúng chỉ dẫn}
\loigiai{
Hành động đúng là “quan sát hình dạng, màu sắc, biểu tượng rồi thực hiện đúng chỉ dẫn” vì giúp nhận biết nhanh loại nguy hiểm, điều cấm và hành động bắt buộc. Chọn D.
}
\end{ex}

% ===== Câu 156 =====
% ID: L10C1B2-58-TN
% Mức: TH
\begin{ex}
% ID: L10C1B2-58-TN
% Mức: TH
Sau khi hoàn thành phép đo liên quan đến biển báo và kí hiệu cảnh báo, hành động nào an toàn nhất?
\choice
{nhờ bạn khác làm thay mà không kiểm tra điều kiện an toàn}
{\True quan sát hình dạng, màu sắc, biểu tượng rồi thực hiện đúng chỉ dẫn}
{bỏ qua biển cảnh báo vì cho rằng chỉ mang tính tham khảo}
{tiếp tục thao tác rồi báo cáo sau}
\loigiai{
Hành động đúng là “quan sát hình dạng, màu sắc, biểu tượng rồi thực hiện đúng chỉ dẫn” vì giúp nhận biết nhanh loại nguy hiểm, điều cấm và hành động bắt buộc. Chọn B.
}
\end{ex}

% ===== Câu 157 =====
% ID: L10C1B2-59-TN
% Mức: TH
\begin{ex}
% ID: L10C1B2-59-TN
% Mức: TH
Khi một bạn chưa đọc hướng dẫn liên quan đến biển báo và kí hiệu cảnh báo, hành động nào an toàn nhất?
\choice
{\True quan sát hình dạng, màu sắc, biểu tượng rồi thực hiện đúng chỉ dẫn}
{bỏ qua biển cảnh báo vì cho rằng chỉ mang tính tham khảo}
{tiếp tục thao tác rồi báo cáo sau}
{nhờ bạn khác làm thay mà không kiểm tra điều kiện an toàn}
\loigiai{
Hành động đúng là “quan sát hình dạng, màu sắc, biểu tượng rồi thực hiện đúng chỉ dẫn” vì giúp nhận biết nhanh loại nguy hiểm, điều cấm và hành động bắt buộc. Chọn A.
}
\end{ex}

% ===== Câu 158 =====
% ID: L10C1B2-62-TN
% Mức: VD
\begin{ex}
% ID: L10C1B2-62-TN
% Mức: VD
Khi có người không đeo phương tiện bảo hộ liên quan đến biển báo và kí hiệu cảnh báo, hành động nào an toàn nhất?
\choice
{nhờ bạn khác làm thay mà không kiểm tra điều kiện an toàn}
{\True quan sát hình dạng, màu sắc, biểu tượng rồi thực hiện đúng chỉ dẫn}
{bỏ qua biển cảnh báo vì cho rằng chỉ mang tính tham khảo}
{tiếp tục thao tác rồi báo cáo sau}
\loigiai{
Hành động đúng là “quan sát hình dạng, màu sắc, biểu tượng rồi thực hiện đúng chỉ dẫn” vì giúp nhận biết nhanh loại nguy hiểm, điều cấm và hành động bắt buộc. Chọn B.
}
\end{ex}

% ===== Câu 159 =====
% ID: L10C1B2-63-TN
% Mức: VD
\begin{ex}
% ID: L10C1B2-63-TN
% Mức: VD
Khi kết quả đo khác dự kiến liên quan đến biển báo và kí hiệu cảnh báo, hành động nào an toàn nhất?
\choice
{\True quan sát hình dạng, màu sắc, biểu tượng rồi thực hiện đúng chỉ dẫn}
{bỏ qua biển cảnh báo vì cho rằng chỉ mang tính tham khảo}
{tiếp tục thao tác rồi báo cáo sau}
{nhờ bạn khác làm thay mà không kiểm tra điều kiện an toàn}
\loigiai{
Hành động đúng là “quan sát hình dạng, màu sắc, biểu tượng rồi thực hiện đúng chỉ dẫn” vì giúp nhận biết nhanh loại nguy hiểm, điều cấm và hành động bắt buộc. Chọn A.
}
\end{ex}

% ===== Câu 160 =====
% ID: L10C1B2-64-TN
% Mức: VD
\begin{ex}
% ID: L10C1B2-64-TN
% Mức: VD
Khi giáo viên yêu cầu dừng thao tác liên quan đến biển báo và kí hiệu cảnh báo, hành động nào an toàn nhất?
\choice
{bỏ qua biển cảnh báo vì cho rằng chỉ mang tính tham khảo}
{tiếp tục thao tác rồi báo cáo sau}
{nhờ bạn khác làm thay mà không kiểm tra điều kiện an toàn}
{\True quan sát hình dạng, màu sắc, biểu tượng rồi thực hiện đúng chỉ dẫn}
\loigiai{
Hành động đúng là “quan sát hình dạng, màu sắc, biểu tượng rồi thực hiện đúng chỉ dẫn” vì giúp nhận biết nhanh loại nguy hiểm, điều cấm và hành động bắt buộc. Chọn D.
}
\end{ex}

